\documentclass[a4paper,11pt]{article}
\usepackage[utf8]{inputenc}
\usepackage[T1]{fontenc}
\usepackage[french]{babel}
\usepackage{geometry}
\geometry{left=1.8cm,right=1.8cm,top=1.3cm,bottom=1.6cm}

\usepackage{lmodern}
\usepackage{xcolor}
\definecolor{titleblue}{RGB}{0,51,140}

\usepackage{hyperref}
\hypersetup{
    colorlinks=true,
    urlcolor=titleblue,
    linkcolor=titleblue,
}

\usepackage{titlesec}
\usepackage{enumitem}

\pagestyle{empty}

% Correction du problème de couleur avec \uppercase
\titleformat{\section}
  {\color{titleblue}\Large\bfseries}{}{0em}{\uppercase}
  [\color{titleblue}\titlerule]

\titlespacing\section{0pt}{8pt}{4pt}
\setlist[itemize]{leftmargin=*,topsep=2pt,itemsep=1pt}

\begin{document}

\begin{center}
    {\Huge\bfseries\color{titleblue} Ismail AISSAOUI}\\[4pt]
    {\Large Ingénieur en Intelligence Artificielle \& Data Science}\\[6pt]
    Email : ismail.aissaoui.pro@gmail.com \quad Tél : +213 660 70 77 96 \\
    Portfolio : \href{https://ismail-aissaoui.vercel.app}{ismail-aissaoui.vercel.app} \quad Localisation : Algérie \\
    LinkedIn : \href{https://linkedin.com/in/aissaoui-ismail-6a77a92a7}{linkedin.com/in/aissaoui-ismail-6a77a92a7}
\end{center}

\section{Résumé}
Ingénieur en IA spécialisé en **Scientific Machine Learning (SciML)** et en **Réseaux de neurones informés par la physique (PINN)**. Expérimenté dans le développement de modèles génératifs pour des systèmes physiques complexes, les dynamiques séquentielles multi-étapes (Mamba-SSM/xLSTM) et la modélisation structurelle via GNN. Déterminé à combler le fossé entre le deep learning et les sciences physiques/chimiques pour résoudre des défis fondamentaux.

\section{Éducation}
\textbf{École Nationale Supérieure d'Informatique (ESI-SBA)} \hfill 2021--2026 \\
Diplôme d'Ingénieur d'État et Master en Informatique \\
Spécialité : **Intelligence Artificielle \& Informatique** \\
*Cours clés : Deep Learning, Apprentissage sur Graphes, Optimisation Métaheuristique, Calcul Scientifique.*

\section{Expertise Technique}
\textbf{Machine Learning} : PyTorch, TensorFlow, Scikit-learn, Lightning, GNN (PyTorch Geometric), ViT. \\
**Physique-IA** : Physics-Informed Neural Networks (PINNs), Modélisation Hybride. \\
**Dynamiques Séquentielles** : Mamba-SSM, xLSTM, Temporal Convolutional Networks (TCN). \\
**Optimisation** : PSO, Algorithmes Génétiques, Optimisation Bayésienne, Optuna. \\
**Déploiement** : FastAPI, Docker, Edge AI (inférence 0,04 ms), CUDA, Linux.

\section{Expérience en Recherche}
\textbf{Thèse d'exercice d'ingénieur (Sonatrach)} \hfill Fév 2026 -- Juil 2026 \\
\textit{Jumeau Numérique Génératif pour Turbines à Gaz GE LM2500}
\begin{itemize}
    \item Développement d'un **Jumeau Numérique informé par la physique** capable de prédire des états thermodynamiques haute-fidélité.
    \item Hybridation des **PINN** avec des architectures séquentielles (**Mamba-SSM** et **xLSTM**) pour capturer des dépendances temporelles non linéaires.
    \item Gestion d'un dataset de **5,2 millions de points (N-CMAPSS)**, optimisation des pipelines de données.
    \item Inférence record de **0,04 ms** sur dispositifs edge via distillation et compilation JIT.
\end{itemize}

\section{Projets Sélectionnés}
\textbf{Aura — Système de Trading IA (LSTM + DQN)} \hfill 2025 \\
Plateforme automatisée utilisant l'apprentissage par renforcement pour optimiser les politiques d'exécution.

\textbf{Elliptic++ — GNN pour la détection de fraude} \hfill 2025 \\
Classification de nœuds illicites sur le réseau Bitcoin via l'apprentissage relationnel.

\textbf{Career Connect — Plateforme RH IA (GNN + RAG)} \hfill 2024 \\
Moteur de matching utilisant les embeddings de graphes et Google Gemini pour l'analyse sémantique.

\section{Langues}
Arabe : Maternel \quad|\quad Français : Courant (C1/C2) \quad|\quad Anglais : Courant (C1/C2)

\end{document}
